\documentclass[11pt,a4paper]{article}

\usepackage[T1]{fontenc}
\usepackage[utf8]{inputenc}
\usepackage{lmodern}
\usepackage{microtype}
\usepackage{geometry}
\geometry{margin=27mm}
\usepackage{amsmath,amssymb,amsthm,mathtools}
\usepackage{enumitem}
\usepackage{booktabs,longtable,array}
\usepackage{xcolor}
\usepackage{hyperref}
\usepackage{url}
\usepackage{fancyhdr}
\usepackage{lastpage}

\hypersetup{
  colorlinks=true,
  linkcolor=blue!55!black,
  citecolor=blue!55!black,
  urlcolor=blue!55!black,
  pdftitle={The Two-Player 3n plus or minus 1 Game - Version 7.0},
  pdfauthor={Grisha Pochuev}
}

\pagestyle{fancy}
\setlength{\headheight}{14pt}
\fancyhf{}
\lhead{Two-Player $3n\pm1$ Game, Version 7.0}
\rhead{Audited status manuscript}
\cfoot{\thepage\ of \pageref{LastPage}}

\newtheorem{theorem}{Theorem}[section]
\newtheorem{proposition}[theorem]{Proposition}
\newtheorem{lemma}[theorem]{Lemma}
\newtheorem{corollary}[theorem]{Corollary}
\theoremstyle{definition}
\newtheorem{definition}[theorem]{Definition}
\newtheorem{conjecture}[theorem]{Conjecture}
\newtheorem{obligation}[theorem]{Open obligation}
\theoremstyle{remark}
\newtheorem{remark}[theorem]{Remark}

\newcommand{\WIN}{\textnormal{\textsc{Win}}}
\newcommand{\LOSS}{\textnormal{\textsc{Loss}}}
\newcommand{\DRAW}{\textnormal{\textsc{Draw}}}
\newcommand{\odd}{\operatorname{odd}}
\newcommand{\Moves}{\operatorname{Moves}}
\newcommand{\asl}{\operatorname{asl}}
\newcommand{\rhoo}{\rho}
\newcommand{\Nat}{\mathbb N}
\newcommand{\Natz}{\mathbb N_0}

\title{\textbf{The Two-Player $3n\pm1$ Game}\\[2mm]
\large An Audited Binary Core, Strengthened Local Normalizations,\\
and the Exact Remaining Global Barrier\\[2mm]
\normalsize Version 7.0}
\author{Grisha Pochuev\\
\texttt{n\_854@mail.ru}\\
\url{https://github.com/Grisha-Pochuev/3n-plus-minus-1-game}}
\date{11 August 2026}

\begin{document}
\maketitle

\noindent
\fbox{\parbox{0.94\textwidth}{
\textbf{Audit status of Version 7.0.}
This document does \emph{not} claim a proof of the all-start no-draw
conjecture.  Earlier versions did make that claim, but hostile review found
that their global marked-router argument used finite states as ranked proof
tokens without first proving the required provenance, and compressed an
attached short-tail module into an unproved exhaustiveness assertion.
Those claims are withdrawn here.  Every unconditional theorem in this
version is proved in the text.  The global conclusion is stated only as a
conditional theorem under four explicit open obligations.  Thus Version 7.0
is intended to be gap-free as a status manuscript, not to conceal an open
step inside a claimed prize solution.
}}

\begin{abstract}
The two-player $3n\pm1$ game is played on positive odd integers.  A move
chooses $3n-1$ or $3n+1$, removes all powers of two, and wins immediately
when the result is $1$.  Arbitrary play can cycle, so the question is whether
optimal play nevertheless has finite remoteness from every start.  We retain
the exact binary conjugation used in earlier manuscripts, in which the two
moves are an expanding map $A$ and a strictly contracting map $B$.  We then
separate four notions that earlier drafts sometimes conflated: a finite
outcome, a routing witness, a retained proof token, and a temporary arithmetic
cursor.

Two local repairs are proved completely.  First, a reverse fixed-tail step
from $Q_r^e(3c)$ either removes one factor of three while preserving an
actual draw, or exposes a genuine loss state and a strictly lower win child
before the draw continues.  Second, the exponent-one predecessor has the
same complete dichotomy; in the exceptional residue classes its two
continuing children form an exact adjacent lift over the actual loss sibling.
We also give a corrected audit of the source-zero valuation-one row and of
the long high-return pair.  These results close several former local gaps,
but they do not establish the global provenance invariant required to
iterate the marked router indefinitely.  The remaining obligations are
isolated precisely, and a complete order-theoretic proof shows that they
would imply the no-draw conjecture.  A companion Python program checks the
arithmetic identities and permanent counterexamples used in the audit; it is
explicitly supporting evidence rather than a proof of the infinite theorem.
\end{abstract}

\tableofcontents

\section{The game, outcomes, and the exact status}

The current state is a positive odd integer $n$.  From $n>1$ the player to
move chooses one of
\[
  3n-1,\qquad 3n+1,
\]
and repeatedly divides the chosen integer by $2$ until it is odd.  Reaching
$1$ wins immediately.  Equivalently, state $1$ is terminal and losing for
the player to move.

For an infinite game graph, the finite-remoteness outcome classes are defined
inductively.  A state is \WIN\ when it has a \LOSS\ child; it is \LOSS\ when
all of its children are \WIN.  These are the least classes closed under the
two rules.  Every remaining state is called \DRAW.

\begin{lemma}[Elementary draw recursion]\label{lem:draw-recursion}
A \DRAW\ state has no \LOSS\ child and has at least one \DRAW\ child.
\end{lemma}

\begin{proof}
A state with a \LOSS\ child is \WIN.  A state all of whose children are
\WIN\ is \LOSS.  Therefore a state in neither finite class has no \LOSS\
child and cannot have all children \WIN; at least one child is \DRAW.
\end{proof}

Arbitrary play does not terminate: in the original game, suitable choices
give $5\to7\to5$.  The target is therefore not termination under all choices.

\begin{conjecture}[No-draw conjecture]\label{conj:main}
Every positive odd starting state has finite remoteness.  Equivalently, the
game has no \DRAW\ positions and optimal play reaches $1$ after finitely
many moves.
\end{conjecture}

Conjecture~\ref{conj:main} is the target of the EUR 500 problem posed by
Ingo Althoefer~\cite{Prize}.  It is \emph{not} promoted to a theorem in this
manuscript.

\section{Binary conjugation}

Write an original odd state as $n=2m+1$ and define
\[
  F(m)=\left\lceil\frac{3m}{2}\right\rceil
       =\frac{3m+(m\bmod2)}{2}.
\]
For later factorization, also put
\[
 J(s)=2F(s)+1=
 \begin{cases}
  3s+1,&s\text{ even},\\
  3s+2,&s\text{ odd}.
 \end{cases}
\]
For a positive integer $x$, let $R(x)$ be the integer obtained by deleting
the maximal alternating suffix of the binary expansion of $x$, and let
$\asl(x)$ denote the number of deleted bits.  Put $R(0)=0$ and
$\asl(0)=0$.
For example,
\[
 R(1001_2)=10_2,
 \qquad
 R(10101_2)=0.
\]

\begin{lemma}[Alternating-suffix factorization]\label{lem:alt-factor}
For every positive integer $z$, put $\delta=1-(z\bmod2)$.  There is an
integer $k\ge1$ such that
\[
  3z+1+\delta=2^k J(R(z)).
\]
If $R(z)>0$, $k$ is the length of the maximal alternating suffix of $z$;
if $R(z)=0$, $k$ is one more than the binary length of $z$.
\end{lemma}

\begin{proof}
First suppose $r=R(z)>0$, let $k=\asl(z)$, and put
$\epsilon=r\bmod2$.  The deleted alternating word has the exact numerical
value
\[
 t_{k,0}=\left\lfloor\frac{2^k}{3}\right\rfloor,
 \qquad
 t_{k,1}=\left\lfloor\frac{2^{k+1}}{3}\right\rfloor,
\]
according as $\epsilon=0$ or $1$, so $z=2^k r+t_{k,\epsilon}$.
Since $J(r)=3r+1+\epsilon$, direct substitution in the two parity cases
gives
\[
  3t_{k,\epsilon}+1+\delta=2^k(1+\epsilon).
\]
Adding the contribution $3\cdot2^k r$ yields the claimed identity.
If $R(z)=0$, the whole binary word is alternating and starts with $1$.
Let $\ell$ be its binary length.  Then $z=t_{\ell,1}$, and the same
calculation gives
\[
  3z+1+\delta=2^{\ell+1}=2^{\ell+1}J(0).
\]
Thus the exponent in the statement is $k=\ell+1$.
\end{proof}

\begin{proposition}[First normal form]\label{prop:first-normal}
For $m>0$, the two children in $m$-coordinates are exactly
\[
  F(m),\qquad F(R(m)).
\]
\end{proposition}

\begin{proof}
Since
\[
  3(2m+1)-1=2(3m+1),\qquad
  3(2m+1)+1=2(3m+2),
\]
exactly one of $3m+1$ and $3m+2$ is odd.  If $m$ is even, the direct odd
expression is $3m+1$ and its new $m$-coordinate is $3m/2=F(m)$.  If $m$
is odd, the direct odd expression is $3m+2$ and its new coordinate is
$(3m+1)/2=F(m)$.

The other expression is
\[
  3m+2-(m\bmod2)=3m+1+\delta,
  \qquad \delta=1-(m\bmod2).
\]
Lemma~\ref{lem:alt-factor} says that its odd part is $J(R(m))$.  Returning
to $m$-coordinates gives
\[
  \frac{J(R(m))-1}{2}=F(R(m)).
\]
Thus the two children are exactly the displayed pair.
\end{proof}

Every child in Proposition~\ref{prop:first-normal} lies in the image of
$F$.  Representing the state $F(q)$ by $q$ gives the conjugated game
\[
  A(q)=F(q)=\left\lceil\frac{3q}{2}\right\rceil,
  \qquad
  B(q)=R(F(q)),
\]
with terminal state $q=0$.

\begin{proposition}[Strict contraction]\label{prop:B-contracts}
For every $q>0$,
\[
  B(q)<q.
\]
\end{proposition}

\begin{proof}
The deletion defining $R$ removes at least one binary digit, so
\[
 B(q)\le \left\lfloor\frac{F(q)}2\right\rfloor
 \le \left\lfloor\frac{3q+1}{4}\right\rfloor<q
\]
for $q\ge2$.  Directly, $B(1)=0$.
\end{proof}

After one original move, the $m$-coordinate is in the image of $F$.
Consequently a proof that every conjugated state $q$ has finite remoteness
would imply Conjecture~\ref{conj:main}; the first original move would then
lead to one of two finite-remoteness states.

\section{Constant tails and coefficient sources}

For a positive odd integer $a$, a phase $e\in\{0,1\}$, and $r\ge1$, define
\[
  Q_r^e(a)=a2^r-e.
\]
This is the unique representation of a positive binary word whose maximal
constant suffix consists of $r$ copies of the bit $e$.

\begin{lemma}[Long-tail recurrence and shared boundary]\label{lem:long-tail}
For every $r\ge3$,
\[
 A(Q_r^e(a))=Q_{r-1}^e(3a),
 \qquad
 B(Q_r^e(a))=Q_{r-2}^e(3a).
\]
Moreover,
\[
 B(Q_1^e(a))=B(Q_2^e(a)).
\]
\end{lemma}

\begin{proof}
The formula for $A$ follows directly from the two phase cases:
\[
 A(a2^r)=3a2^{r-1},
 \qquad
 A(a2^r-1)=3a2^{r-1}-1.
\]
When $r\ge3$, the image still ends in at least two equal bits.  Its maximal
alternating suffix therefore consists only of the last bit, and applying
$R$ removes one further constant bit, proving the formula for $B$.

At the boundary,
\[
 A(Q_1^e(a))=3a-e,
 \qquad
 A(Q_2^e(a))=6a-e=2(3a-e)+e.
\]
The last bit of $3a-e$ is $1-e$, so the appended bit $e$ extends the same
maximal alternating suffix by one digit.  Deleting it leaves the same
prefix in both cases.
\end{proof}

Recall that $J(s)=2A(s)+1$.  The map $J$ is an increasing bijection
from $\Natz$ to the positive odd integers not divisible by $3$: explicitly,
$6u+1=J(2u)$ and $6u+5=J(2u+1)$.  Thus every positive odd coefficient has a
unique factorization
\[
